\chapter{Model and Regularisation Methods}
\label{chap:main-methods}

Functional linear regression relates a scalar response to a predictor curve.
The difficulty is to recover the curve's effect when neighbouring measurements
are strongly related and some directions contain little predictor variation.
The four procedures studied here address this difficulty by restricting the
directions used for estimation. They differ in how those directions are
constructed, how the fit is calculated, and how its complexity is selected.
This chapter gives the definitions needed to interpret the results. Extended
background, proofs, and algorithmic details are retained in
Appendices~\ref{chap:background}, \ref{app:mathematical-proofs}, and
\ref{chap:methods}, respectively.

\section{The model and the identified slope}

Let $X$ be a random element of the real Hilbert space
$\mathcal H=L^2([0,1])$, with inner product
$\langle f,g\rangle=\int_0^1 f(t)g(t)\,\mathrm dt$ and
$\mathbb E\|X\|^2<\infty$. Measurements on another compact interval can be
rescaled to $[0,1]$. The model is
\begin{equation}
 Y=\alpha+\langle X,\beta\rangle+\varepsilon,
 \qquad \mathbb E(\varepsilon\mid X)=0,
 \qquad \beta\in\mathcal H.
 \label{main:eq:model}
\end{equation}
Here $Y$ has finite second moment. Writing $X^c=X-\mathbb EX$ and
$Y^c=Y-\mathbb EY$ gives
$Y^c=\langle X^c,\beta\rangle+\varepsilon$; the intercept is recovered as
$\alpha=\mathbb EY-\langle\mathbb EX,\beta\rangle$. Treating the
predictor as a function preserves the ordering and proximity of its
measurement locations, even though computation ultimately uses a finite
grid or basis \parencite{ramsay_silverman_2005,hsing_eubank_2015}.

Define the covariance operator and cross-covariance element by
\begin{equation}
 Kf=\mathbb E\{\langle X^c,f\rangle X^c\},
 \qquad r=\mathbb E(Y^cX^c),
 \qquad K\beta=r.
 \label{main:eq:normal-equation}
\end{equation}
The last equality follows from the conditional mean restriction. The
operator $K$ is self-adjoint, non-negative, and trace class. Its positive
eigenvalues $\lambda_j$ and orthonormal eigenfunctions $v_j$ describe the
directions and amounts of predictor variation. If $K$ has infinite rank,
$\lambda_j\to0$.

Two slopes that differ by an element of $\operatorname{Null}(K)$ give the
same centred conditional mean. The data therefore identify only their
equivalence class. Throughout, $\beta$ denotes its minimum-norm
representative in
$\mathcal H_{\mathrm{id}}=\operatorname{Null}(K)^\perp
=\overline{\operatorname{Range}(K)}$.
On this space, the formal inverse has coefficients
$\beta_j=\langle r,v_j\rangle/\lambda_j$. Small eigenvalues amplify errors
in the estimated moments. Consequently, identification does not guarantee
stable estimation: even an injective covariance operator can have an
unbounded inverse. Regularisation limits the contribution of directions in
which this amplification is severe
\parencite{engl_hanke_neubauer_1996,hall_horowitz_2007}.

The spectral expansion also clarifies what is being estimated. The scores
$\xi_j=\langle X^c,v_j\rangle$ are uncorrelated, with variance
$\lambda_j$, and $X^c=\sum_j\xi_jv_j$ in mean square. Its contribution
to the response is $\sum_j\xi_j\beta_j$. A principal direction can
therefore explain much of the predictor variation yet little of the
response if its slope coefficient is small. Conversely, a lower-variance
direction can matter for prediction when its slope coefficient is large.
The simulation designs change the covariance spectrum and slope
coefficients separately to examine these possibilities.

\section{Slope recovery and prediction}

The relevant loss depends on the statistical question. For
$R(b)=\mathbb E\{Y^c-\langle X^c,b\rangle\}^2$, the model implies
\begin{equation}
 R(b)-R(\beta)
 =\langle K(b-\beta),b-\beta\rangle
 =\|K^{1/2}(b-\beta)\|^2.
 \label{main:eq:prediction-risk}
\end{equation}
Slope recovery instead uses the integrated squared error $\|b-\beta\|^2$.
A large slope error in a low-variance direction may have little effect on
prediction. The covariance-weighted expression in
\eqref{main:eq:prediction-risk} is a squared seminorm on $\mathcal H$ and
a squared norm on $\mathcal H_{\mathrm{id}}$; it generally controls less
than the ordinary squared norm. This explains why accurate predictions do
not establish accurate slope recovery \parencite{cai_hall_2006}.

For example, if the estimation error is $av_j$ along one covariance
eigenfunction, its integrated squared error is $a^2$, whereas its excess
prediction risk is $\lambda_j a^2$. As $\lambda_j$ becomes small, the
same slope error becomes less visible in prediction. This comparison
concerns excess risk above the irreducible response noise. An independent
test mean squared prediction error also includes that noise, so the
simulation reports both slope loss and prediction loss instead of treating
either as a complete assessment.

Regularisation balances the bias introduced by suppressing weak directions
against the variability incurred by estimating them. Spectral truncation
discards selected directions explicitly. Iterative methods restrict a
sequence of increasingly rich spaces, making the stopping index a
regularisation parameter. Continuing until the empirical equation is solved
as accurately as possible can fit sampling noise. Inference introduces a
further requirement on the remaining moment residual, considered in
Chapter~\ref{chap:main-inference}.

Other regularisation families provide useful context. A roughness penalty
controls derivatives of the fitted slope, while Tikhonov regularisation
replaces inverse weights $1/\lambda$ by $1/(\lambda+\rho)$ for
$\rho>0$ on the positive spectrum. Both can stabilise weak directions
without selecting a component count
\parencite{cardot_ferraty_sarda_2003,crambes_kneip_sarda_2009}.
The present comparison concentrates on spectral and Krylov procedures
because they allow the roles of subspace choice, numerical representation,
and stopping to be examined together. Its conclusions concern these four
complete procedures, rather than an exhaustive ranking of all functional
regression estimators.

\section{The empirical problem}

Let $x_i$ and $y_i$ denote the centred training curves and responses. The
simulations use their known zero population means; the empirical study
estimates means within each training sample. Define
$(T_nb)_i=\langle x_i,b\rangle$ and equip response space with
$\|a\|_n^2=n^{-1}a^\top a$. Then
$T_n^*a=n^{-1}\sum_i a_ix_i$,
$\widehat K=T_n^*T_n$, and $\widehat r=T_n^*y$.
The response residual is $y-T_nb$, whereas the moment residual is
$e(b)=\widehat r-\widehat Kb=T_n^*(y-T_nb)$.
These residuals belong to different spaces; minimising their norms need not
give the same restricted estimate.

On the equally weighted $T$-point grid, the functional inner product becomes
$\langle f,g\rangle_T=T^{-1}f^\top g$. For the $n\times T$ matrix $X$ of
centred curves, this gives
\begin{equation}
 \widehat K=\frac{X^\top X}{nT},\qquad
 \widehat r=\frac{X^\top y}{n},\qquad
 T_nb=\frac{Xb}{T}.
 \label{main:eq:grid-moments}
\end{equation}
The scaling is shared by all four methods. With sample centring,
$\operatorname{rank}(\widehat K)\leq n-1$; known-mean centring gives the
bound $n$. Numerical rank safeguards further restrict usable directions.
Discretisation makes the calculation finite, but does not eliminate
ill-conditioning or make a fixed grid identical to its continuous analogue.
The estimation comparison uses the common candidate range
$m=1,\ldots,70$, subject to each method's rank and validity rules; a
zero-component estimator is not included. The selected point on the path
remains method-specific. The more demanding stopping calculations used for
inference are described separately in Chapter~\ref{chap:main-inference}.

\section{Functional principal component regression}

FPCR retains the leading empirical covariance eigenfunctions. For
$\mathcal H$-normalised eigenpairs
$(\widehat\lambda_j,\widehat v_j)$, its $m$-component slope is
\begin{equation}
 \widehat\beta_m^{\mathrm{FPCR}}
 =\sum_{j=1}^{m}
 \frac{\langle\widehat r,\widehat v_j\rangle}
      {\widehat\lambda_j}\widehat v_j.
 \label{main:eq:fpcr}
\end{equation}
This is response least squares over the first $m$ principal directions.
Only numerically positive eigenvalues are eligible. FPCR uses predictor
variance to order directions, independently of their relationship with the
response. It is consequently well suited to signals concentrated in leading
eigenfunctions, but a response-relevant direction with modest predictor
variance can enter late \parencite{cardot_ferraty_sarda_1999}.

The reported FPCR fit selects the smallest minimiser of response-space
generalised cross-validation,
\begin{equation}
 \operatorname{GCV}_{Y}(m)
 =\frac{\|y-T_n\widehat\beta_m^{\mathrm{FPCR}}\|_n^2}
             {(1-m/n)^2}.
 \label{main:eq:response-gcv}
\end{equation}
This selection rule is distinct from the internal FPCR pilot used to
estimate the noise scale for CG stopping below.

For a fixed retained eigenspace, the fitted responses are an orthogonal
projection onto $m$ score directions, which motivates the dimension
penalty in \eqref{main:eq:response-gcv}. This interpretation concerns the
fit at a fixed $m$; selection itself also depends on the observations.
Adding a principal component always enlarges the available response
space, but may admit a poorly estimated slope coefficient through division
by a small eigenvalue. The best in-sample response fit is consequently not
the criterion for choosing the final dimension.

\section{Raw and Arnoldi FPLS}

Functional partial least squares uses the response when constructing its
directions. Its alternative PLS representation searches the Krylov space
\parencite{delaigle_hall_2012}
\begin{equation}
 \begin{aligned}
 \widehat{\mathcal K}_m
 &=\operatorname{span}\{\widehat r,\widehat K\widehat r,
       \ldots,\widehat K^{m-1}\widehat r\},\\
 \widehat\beta_m^{\mathrm{APLS}}
 &\in\arg\min_{b\in\widehat{\mathcal K}_m}\|y-T_nb\|_n^2.
 \end{aligned}
 \label{main:eq:apls}
\end{equation}
The cross-covariance $\widehat r$ makes this space response-informed.
Classical NIPALS-type score construction and deflation yield the same
fixed-$m$ response fit under the stated full-rank conditions; NIPALS is a
reference formulation, not a fifth comparator. The corresponding
first-order condition is
$e(\widehat\beta_m^{\mathrm{APLS}})\perp\widehat{\mathcal K}_m$.

Raw FPLS forms the power-basis matrix
$G_m=[\widehat r,\widehat K\widehat r,\ldots,
\widehat K^{m-1}\widehat r]$ directly, without column rescaling or
orthogonalisation. With $Z_m=T_nG_m$, it solves
$(Z_m^\top Z_m)a=Z_m^\top y$ and returns $G_ma$.
Repeated covariance powers increasingly emphasise the largest active
eigenvalues, making the columns nearly parallel. Forming normal equations
then squares the design condition number whenever $Z_m$ has full column
rank. A failed or non-finite solve is marked invalid; it is not silently
replaced by a pseudoinverse or an orthogonalised method.

More explicitly, in the empirical eigenbasis,
$\widehat K^k\widehat r=\sum_j\widehat\lambda_j^k
\langle\widehat r,\widehat v_j\rangle\widehat v_j$.
Relative contributions shrink like
$(\widehat\lambda_j/\widehat\lambda_\star)^k$, where
$\widehat\lambda_\star$ is the largest eigenvalue with a nonzero
cross-covariance projection. Distinct raw columns can then carry almost
the same numerical information. Representing a useful direction by
cancelling large, nearly equal contributions makes the coefficient solve
sensitive to rounding. An algebraically correct span alone does not
guarantee that its power-basis coordinates can be used reliably.

Arnoldi FPLS changes the numerical basis while retaining
\eqref{main:eq:apls}. Starting from the normalised cross-covariance, it
applies $\widehat K$ successively and removes projections onto the existing
basis twice, using two-pass classical Gram--Schmidt. The resulting
orthonormal matrix $Q_m$ represents the same Krylov space in exact
arithmetic. Response least squares on $T_nQ_m$ is solved by QR
factorisation. Fixed rank tests terminate basis expansion or freeze the
response path at its last valid fit when additional directions are
numerically unavailable. The precise thresholds and failure rules are
documented in Appendix~\ref{chap:methods}.

In grid coordinates, the columns of $Q_j$ are Euclidean-normalised, so
$Q_j^\top Q_j=I$. One Arnoldi step starts with $z=\widehat Kq_j$,
applies $z\leftarrow z-Q_j(Q_j^\top z)$ twice, and then divides by
$\|z\|_2$. The Euclidean and discrete functional norms differ only by
the common factor $\sqrt T$; this constant scaling changes neither the
Krylov span nor which directions are orthogonal. Repeating the
projection reduces loss of orthogonality, while QR avoids squaring the
response-design condition number. These operations stabilise how the
restricted problem is represented; they do not remove the statistical
uncertainty associated with adding weak directions. Stopping or component
selection is still needed.

At a common valid dimension, Raw and Arnoldi therefore share the same
exact-arithmetic target. This does not force their selected estimates to
coincide. Each method minimises its own five-fold prediction-error curve,
using identical folds. A raw dimension is eligible only if its full-sample
fit and every training-fold fit are valid and finite. Numerical instability
can change both its eligible dimensions and its cross-validation curve.
The smallest minimiser is selected, and the chosen procedure is refitted
on the full training sample.

\section{CG--FPLS and discrepancy stopping}

The CG--FPLS procedure follows \textcite{babii_carrasco_tsafack_2026}, but
the label requires a precise interpretation. Its finite-iteration target is
\begin{equation}
 \widehat\beta_m^{\mathrm{CG}}
 \in\arg\min_{b\in\widehat{\mathcal K}_m}
       \|\widehat r-\widehat Kb\|^2,
 \qquad e_m=\widehat r-\widehat K\widehat\beta_m^{\mathrm{CG}}.
 \label{main:eq:cg}
\end{equation}
It searches the same nested spaces as APLS, but minimises the moment
residual. Its first-order condition is
$e_m\perp\widehat K\widehat{\mathcal K}_m$, rather than
$e_m\perp\widehat{\mathcal K}_m$. Thus it generally gives a different
finite-$m$ slope from response-space FPLS.

The estimation implementation uses a short recurrence to avoid explicitly
forming the power-basis equations. In standard numerical-linear-algebra
terminology, these residual-minimising updates are the
\emph{conjugate-residual} recurrence, rather than textbook energy-norm
conjugate gradient. The thesis retains the CG--FPLS name while distinguishing
the actual objective and updates. Algebraic equivalences refer to valid
iterates in exact arithmetic; the numerical breakdown rules are specified
in Appendix~\ref{chap:methods}.

For completeness, initialise $\widehat\beta_0=0$ and
$e_0=d_0=\widehat r$. With $u_k=\widehat Kd_k$, the updates are
\begin{align}
 \alpha_k&=\frac{\langle e_k,\widehat K e_k\rangle}{\|u_k\|^2},
 \label{main:eq:cr-alpha}\\
 \widehat\beta_{k+1}&=\widehat\beta_k+\alpha_kd_k,
 \label{main:eq:cr-slope}\\
 e_{k+1}&=e_k-\alpha_ku_k,
 \label{main:eq:cr-residual}\\
 \gamma_k&=\frac{\langle e_{k+1},\widehat K e_{k+1}\rangle}
                    {\langle e_k,\widehat K e_k\rangle},
 \label{main:eq:cr-gamma}\\
 d_{k+1}&=e_{k+1}+\gamma_kd_k.
 \label{main:eq:cr-direction}
\end{align}
In exact arithmetic the images of the search directions under
$\widehat K$ are mutually orthogonal while the recurrence remains valid.
The new direction therefore permits further moment-residual reduction
without undoing the earlier minimisations. This is the relevant
conjugacy for the objective in \eqref{main:eq:cg}. The implementation
symmetrises the empirical operator and guards against non-positive or
non-finite quadratic forms and negligible direction images. Failure of
these safeguards freezes the remaining path at its last valid estimate;
non-finite updated iterates or recurrence coefficients raise an error.
These numerical safeguards should be distinguished from satisfying the
statistical stopping threshold.

The estimation rule selects the first $m\geq1$ such that
\begin{equation}
 \|e_m\|_T\leq\widehat\nu_n,
 \qquad
 \widehat\nu_n=1.01\left(
       \frac{2\widehat\sigma^2\widehat E_X}{0.1n}
                    \right)^{1/2},
 \qquad \widehat E_X=\frac1n\sum_i\|x_i\|_T^2.
 \label{main:eq:cg-threshold}
\end{equation}
The noise estimate $\widehat\sigma^2$ is the mean squared response
residual of an internal FPCR pilot. Its dimension minimises moment-space
GCV, $\|\widehat r-\widehat K\widehat\beta_q^{\mathrm{FPCR}}\|_T^2/
(1-q/T)^2$, over the permitted pilot path. The estimation studies allow
up to 70 pilot components, subject to numerical rank. This pilot supplies
the threshold scale; it does not add another reported method. If the
threshold is not reached, the maximum permitted component count is used
and the absence of a crossing is recorded.

The pilot matters because it determines how small a residual is regarded
as sufficient. Holding the path and predictor energy fixed, a smaller
estimated noise variance lowers the threshold and can delay stopping;
a larger estimate can stop the procedure earlier. The pilot's own
moment-GCV criterion is therefore part of the CG procedure's behaviour,
even though the pilot slope is not reported as the final estimate.

Discrepancy stopping compares the residual with an estimated sampling-noise
scale. The associated ideal rule uses population noise variance and
predictor energy. Its theoretical guarantees, under the required moment
and source conditions, do not automatically establish those of this
particular plug-in pilot. The empirical performance of the complete
procedure must therefore be assessed directly.

In particular, a smaller residual is not by itself evidence of a better
slope estimate. The residual is measured after applying the covariance
operator, and therefore discounts errors in weak directions even more
strongly than prediction risk. Later steps may continue reducing this
residual while increasing sensitivity to sampling error. The source
condition used in the theoretical analysis has the form $\beta=K^\mu g$,
which expresses how strongly the slope is concentrated in well-identified
directions. Unknown slope regularity motivates adaptive stopping, but does
not make the same stopping choice suitable for estimation and inference
\parencite{babii_carrasco_tsafack_2026}.

\section{What the four-method comparison identifies}

Table~\ref{main:tab:methods} summarises the three choices that define a
procedure: its search space, restricted objective, and selection rule.
The table also explains why a component count cannot be interpreted
independently of the method producing it.

\begin{table}[htbp]
 \centering
 \small
 \caption[The four complete estimation procedures]{Search spaces, finite-component
 objectives, and tuning rules. The maximum estimation path is 70 components,
 subject to numerical rank and validity safeguards.}
 \label{main:tab:methods}
 \begin{tabular}{@{}
 >{\raggedright\arraybackslash}p{1.8cm}
 >{\raggedright\arraybackslash}p{3.0cm}
 >{\raggedright\arraybackslash}p{3.2cm}
 >{\raggedright\arraybackslash}p{4.4cm}@{}}
 \toprule
 Method & Search space & Objective & Dimension selection \\
 \midrule
 FPCR & Leading covariance eigenfunctions & Response least squares
 & Response GCV \\
 Raw FPLS & Krylov space in the raw power basis & Response least squares
 & Five-fold CV over valid dimensions \\
 Arnoldi FPLS & Same Krylov space in an orthonormal basis & Response least squares
 & Five-fold CV, separately minimised on the same folds \\
 CG--FPLS & Same nested Krylov spaces & Moment-residual least squares
 & Discrepancy stopping with an internal noise pilot \\
 \bottomrule
 \end{tabular}
\end{table}

The comparisons that follow retain each method's own objective and tuning
rule. A common path cap does not imply equal effective complexity, and a
selected component count is not a common measure of degrees of freedom
across all four methods. Raw versus Arnoldi primarily exposes numerical
representation; Arnoldi versus CG also changes the restricted objective
and stopping rule. These distinctions guide the interpretation of both the
simulation results and the spectroscopy application.
